Al jaciment neolític de la Draga (Banyoles) s'utilitza el mètode del carboni 14 (${}^{14}\mathrm{C}$) per a la datació de restes. El ${}^{14}\mathrm{C}$ és un isòtop radioactiu que es produeix a l'atmosfera de manera continuada quan neutrons provinents dels raigs còsmics xoquen amb els nuclis de nitrogen (${}^{14}\mathrm{N}$). Els éssers vius l'incorporem a l'organisme mentre interactuem amb l'atmosfera, però, a partir del moment de la mort, el ${}^{14}\mathrm{C}$ comença a desintegrar-se novament en ${}^{14}\mathrm{N}$, amb un període de semidesintegració de 5\,730 anys. La datació d'una mostra d'os d'animal del sector D de la Draga és de 6\,010 anys.

\begin{apartat}{1,25}
Escriviu la reacció nuclear que correspon al decaïment del ${}^{14}\mathrm{C}$ a ${}^{14}\mathrm{N}$, incloent-hi els antineutrins. Justifiqueu de quin tipus de reacció nuclear es tracta. A partir de l'equació de la desintegració, determineu la relació entre la constant de desintegració $\lambda$ i el període de semidesintegració.
\end{apartat}

\begin{apartat}{1,25}
Calculeu la constant de desintegració $\lambda$ del ${}^{14}\mathrm{C}$. Determineu quin percentatge de ${}^{14}\mathrm{C}$ encara roman a la mostra d'os del sector D respecte al ${}^{14}\mathrm{C}$ que hi havia quan l'organisme estava viu. Calculeu l'activitat de l'os en el moment de morir l'animal si actualment l'activitat de l'os és de 3,63~Bq.
\end{apartat}

\dades{%
  El nombre atòmic del carboni és $Z = 6$.\\
  El nombre atòmic del nitrogen és $Z = 7$.}
